% Adapted from the author's public Erdos Paper IV v1.23 series_template.tex.
% Layout provenance is recorded in 02_validation; the Markdown is semantic authority.
\documentclass[spanish,11pt,a4paper]{article}
\usepackage[margin=1in]{geometry}
\usepackage{fontspec}
\setmainfont{lmroman10-regular.otf}[BoldFont=lmroman10-bold.otf,ItalicFont=lmroman10-italic.otf,BoldItalicFont=lmroman10-bolditalic.otf]
\setsansfont{lmsans10-regular.otf}[BoldFont=lmsans10-bold.otf]
\setmonofont{lmmono10-regular.otf}[Scale=MatchLowercase]
\usepackage{amsmath,amssymb}
\providecommand{\mathscr}{\mathcal}
\usepackage[spanish,es-noshorthands,es-tabla]{babel}
\usepackage{microtype,parskip}
\usepackage{longtable,booktabs,array,calc}
\usepackage{xcolor,xurl}
\usepackage{hyperref,bookmark}
\usepackage{etoolbox,needspace}
\setcounter{secnumdepth}{-\maxdimen}
\setlength{\emergencystretch}{3em}
\providecommand{\tightlist}{\setlength{\itemsep}{0pt}\setlength{\parskip}{0pt}}
\newcounter{none}
\AtBeginEnvironment{longtable}{\Needspace{8\baselineskip}\small}
\pretocmd{\section}{\Needspace{5\baselineskip}}{}{}
\pretocmd{\subsection}{\Needspace{7\baselineskip}}{}{}
\pretocmd{\subsubsection}{\Needspace{6\baselineskip}}{}{}
\urlstyle{same}

\usepackage{graphicx}
\providecommand{\pandocbounded}[1]{#1}

\usepackage{float}

\newcommand{\KL}{\operatorname{KL}}
\clubpenalty=10000\widowpenalty=10000\displaywidowpenalty=10000
\hypersetup{pdftitle={Documento formal complementario: entropía y grandes desviaciones para cuadrados latinos},pdfauthor={Juan Pablo Traverso Gianini},pdflang={es},colorlinks=true,linkcolor=blue,urlcolor=blue,citecolor=blue}
\title{Documento formal complementario: entropía y grandes desviaciones para cuadrados latinos}
\author{Juan Pablo Traverso Gianini\\\small Independent researcher, Santiago, Chile\\\small \href{mailto:jtraverso@gmail.com}{jtraverso@gmail.com}\\\small \href{https://orcid.org/0009-0003-6068-4096}{ORCID: 0009-0003-6068-4096}}
\date{}\begin{document}\raggedbottom\maketitle
\textbf{Preprint 1.0 --- 16 de noviembre de 2026.}

\textbf{Repositorio:} \url{https://github.com/jtraverso/latinon-entropy-ldp}.

\subsection{1. Alcance e interfaz matemática}\label{alcance-e-interfaz-matemuxe1tica}

Este documento registra cómo se representan en Lean los enunciados del manuscrito. Es una guía de las fuentes y evidencias entregadas, no un reemplazo de las demostraciones escritas. El modelo tiene tres márgenes uniformes de dos coordenadas y conserva el mapa de orden que preserva la medida. Su distancia tiene dos términos de orden y un término de distribución que prueba intervalos de símbolos. La entropía relativa toma valores en los reales no negativos extendidos. Su conversión a un número real solo se justifica después de una hipótesis de tasa finita.

Los enunciados formales son los contratos protegidos del registro de auditoría. El resultado superior usa \texttt{ENNReal.ofReal\ t}; este es el piso no estricto max(t,0) del manuscrito. El resultado inferior requiere un costo testigo finito no negativo. Ninguno presupone una ley sobre una colección auxiliar: ambos se refieren al transporte de la ley uniforme finita real de cuadrados latinos.

\subsection{2. Correspondencia de enunciados principales}\label{correspondencia-de-enunciados-principales}

\Needspace{7\baselineskip}\textbf{Theorem 2.1 / Teorema 2.1}

\path{LatinonEntropy.closed_uniform_latinon_upper}

Fuente: \path{lean/LatinonBridge/LatinonGlobalDeviations.lean}, línea 20.

\path{LatinonEntropy.open_uniform_latinon_lower}

Fuente: \path{lean/LatinonBridge/LatinonGlobalDeviations.lean}, línea 46.

\Needspace{7\baselineskip}\textbf{Theorem 2.2 / Teorema 2.2}

\path{LatinonEntropy.finite_rate_standard_pattern_log_count}

Fuente: \path{lean/LatinonPatternConjecture.lean}, línea 12.

\path{LatinonEntropy.empty_standard_pattern_fiber_eventually_zero}

Fuente: \path{lean/LatinonPatternConjecture.lean}, línea 42.

\Needspace{7\baselineskip}\textbf{Theorem 3.1 / Teorema 3.1}

\path{LatinonEntropy.finite_latin_tensor_lower_bound}

Fuente: \path{lean/LatinTensorLowerBound.lean}, línea 51.

\Needspace{7\baselineskip}\textbf{Lemma 4.1 / Lema 4.1}

\path{LatinonEntropy.correction_lemma}

Fuente: \path{lean/QuasirandomCorrection.lean}, línea 143.

La proposición 2.3 es la identificación escrita con las representaciones y la definición de distancia de corte publicadas por GHHS. Su prueba mediante desintegración y cambio de variables está en el manuscrito; no se presenta como una declaración Lean nueva. La fuente formal usa la representación por medidas, transportes inversos de coordenadas y dos términos de orden. La notación del manuscrito no renombra ninguna declaración Lean.

\textbf{Tasa buena en el teorema 2.1}

\path{LatinonEntropy.latinon_rate_lowerSemicontinuous}

Fuente: \path{lean/LatinonFullDyadicKL.lean}, línea 42.

\path{LatinonEntropy.latinon_full_cut_compact}

Fuente: \path{lean/LatinonBridge/LatinonCompactness.lean}, línea 60.

\path{LatinonEntropy.latinon_rate_sublevels_compact}

Fuente: \path{lean/LatinonBridge/LatinonGlobalDeviations.lean}, línea 15.

Estas tres propiedades se consultan junto con los ocho resultados originales y el teorema de corrección reforzado en \path{ReviewAxioms.lean}. Las doce consultas actuales informan únicamente axiomas fundacionales estándar. Las once salidas originales también se conservan con la reconstrucción limpia en \path{preprints/LATINON/02_validation/CLEAN_REBUILD_AUDIT.json}.

El teorema finito de tensores cuantifica el umbral de orden antes del perfil. El teorema de corrección no tiene hipótesis de diámetro residual ni certificados de potenciales. El lema 4.1 conserva sus constantes explícitas de caja y el coeficiente lineal 3.

El argumento de diámetro seis también aparece internamente como \path{LatinonEntropy.ScaledResidualData.diameter_six} en \path{lean/ResidualDiameter.lean}. Su ausencia de las hipótesis de \path{correction_lemma} significa que se deduce internamente, no que solo exista en papel. Se revisa el lema 4.1 como prueba escrita por separado de su correspondencia con ese argumento de fuentes.

\textbf{Mejora del lema 4.1.}

\path{LatinonEntropy.ScaledResidualData.diameter_five}

\path{LatinonEntropy.correction_lemma_five}

Ambas declaraciones están en \path{lean/CorrectionFive.lean}. La transposición invierte las aristas residuales; sus caminos de cuatro pasos entre filas dan caminos de cuatro pasos entre columnas en el grafo original. Los caminos entre lados opuestos usan a lo sumo cinco aristas, lo que da el radio del potencial y el coeficiente de corrección \(5/2\). Estas declaraciones refuerzan el resultado original de coeficiente tres sin modificar sus consumidores ni la correspondencia de teoremas numerados.

\subsection{3. Registro de dependencias}\label{registro-de-dependencias}

El registro central tiene 23 pasos. La lista siguiente da la declaración y las dependencias en cada paso. La correspondencia actual con las fuentes compiladas está en \path{publication/PROOF_LEDGER.json}. Los certificados originales ligados a fuentes que se conservan dentro del archivo Lean son evidencia histórica; el informe de reconstrucción completada certifica la selección entregada.

\Needspace{6\baselineskip}\textbf{1. Modelo y distancia completa.}

\path{LatinonEntropy.LatinonMeasure}

Dependencias en el registro: ninguna.

\Needspace{6\baselineskip}\textbf{2. Corrección determinista.}

\path{LatinonEntropy.correction_lemma}

Dependencias en el registro: ninguna.

\Needspace{6\baselineskip}\textbf{3. Permanentes y marginales.}

\path{LatinonEntropy.permanental_marginal_l1_le}

Dependencias en el registro: ninguna.

\Needspace{6\baselineskip}\textbf{4. Proceso latino y parada.}

\path{LatinonEntropy.profileProcessPolicy}

Dependencias en el registro: 2, 3.

\Needspace{6\baselineskip}\textbf{5. Recurrencia del defecto.}

\path{LatinonEntropy.profile_expected_defect_recurrence}

Dependencias en el registro: 4.

\Needspace{6\baselineskip}\textbf{6. Estabilidad por Gronwall.}

\path{LatinonEntropy.profile_expected_defect_uniform}

Dependencias en el registro: 5.

\Needspace{6\baselineskip}\textbf{7. Supervivencia uniforme.}

\path{LatinonEntropy.eventually_uniform_profile_survival}

Dependencias en el registro: 6.

\Needspace{6\baselineskip}\textbf{8. Cercanía uniforme en cortes.}

\path{LatinonEntropy.eventually_uniform_profile_closeness}

Dependencias en el registro: 7.

\Needspace{6\baselineskip}\textbf{9. Entropía del proceso completo.}

\path{LatinonEntropy.profile_history_entropy_lower}

Dependencias en el registro: 6, 7.

\Needspace{6\baselineskip}\textbf{10. Conteo inferior de tensores.}

\path{LatinonEntropy.finite_latin_tensor_lower_bound}

Dependencias en el registro: 8, 9.

\Needspace{6\baselineskip}\textbf{11. Puente tensor-Latinon.}

\path{LatinonEntropy.finite_latin_grid_metric_ball_lower_bound}

Dependencias en el registro: 1, 10.

\Needspace{6\baselineskip}\textbf{12. Cota superior de particiones.}

\path{LatinonEntropy.exists_latinon_partition_ball_count_upper}

Dependencias en el registro: 1.

\Needspace{6\baselineskip}\textbf{13. Tasa KL extendida.}

\path{LatinonEntropy.latinon_rate_eq_iSup_dyadic}

Dependencias en el registro: 12.

\Needspace{6\baselineskip}\textbf{14. Centro fijo y orden.}

\path{LatinonEntropy.exists_bounded_latinon_profile_approximation}

Dependencias en el registro: 11, 13.

\Needspace{6\baselineskip}\textbf{15. Regularización y entropía local.}

\path{LatinonEntropy.finite_rate_latinon_entropy_count}

Dependencias en el registro: 14.

\Needspace{6\baselineskip}\textbf{16. Compacidad de la distancia completa.}

\path{LatinonEntropy.latinon_full_cut_compact}

Dependencias en el registro: 1.

\Needspace{6\baselineskip}\textbf{17. LDP global.}

\path{LatinonEntropy.closed_uniform_latinon_upper}

Dependencias en el registro: 13, 15, 16.

\Needspace{6\baselineskip}\textbf{18. Contracción y observables.}

\path{LatinonEntropy.latinonContractedRate_attained}

Dependencias en el registro: 17.

\Needspace{6\baselineskip}\textbf{19. Densidades,continuidad y normalización.}

\path{LatinonEntropy.sum_latinonPatternDensity}

Dependencias en el registro: 1.

\Needspace{6\baselineskip}\textbf{20. Puente finito de coordenadas.}

\path{LatinonEntropy.finite_matrix_pattern_embedding_error}

Dependencias en el registro: 19.

\Needspace{6\baselineskip}\textbf{21. Normalización binomial.}

\path{LatinonEntropy.standard_pattern_profile_uniform_error}

Dependencias en el registro: 20.

\Needspace{6\baselineskip}\textbf{22. Conteo con restricciones cerradas.}

\path{LatinonEntropy.finite_rate_standard_pattern_entropy_count}

Dependencias en el registro: 18, 21.

\Needspace{6\baselineskip}\textbf{23. Fórmula logarítmica y fibra vacía.}

\path{LatinonEntropy.finite_rate_standard_pattern_log_count}

Dependencias en el registro: 22, 16.

\subsection{4. Aplicaciones y ensamblajes de fuentes}\label{aplicaciones-y-ensamblajes-de-fuentes}

El ensamblaje de aplicaciones es \path{LatinonSeriesApplications.lean}. Contiene las definiciones CDR finitas y el argumento de huella de pares de Wilson, seguido de la correspondencia de factorizaciones, la sustitución de packings, los productos de entropía local y la conservación exacta de tipos. Sus ocho entradas originales se conservan en \path{lean/evidence/applications}. El ensamblaje central de patrones conserva de forma análoga sus 18 entradas en \path{lean/evidence/proof-inputs}. Conservar una entrada no afirma que se compile como módulo independiente: los objetivos de compilación son los módulos ensamblados.

\textbf{A1.} \path{finite_rate_factorization_entropy_count}

\textbf{A2.} \path{finite_rate_factorization_pattern_log_count}

\textbf{A3.} \path{replaceSteinerPacking_bicoloured}

\textbf{A4.} \path{wilsonLatinBlockTemplate}

\textbf{A5.} \path{card_restrictedRealizations}

\textbf{A6.} \path{wilson_distinguished_latinon_blocks_entropy_count}

\textbf{A7.} \path{wilson_completed_exact_type_count_eq_base}

\textbf{A8.} \path{residual_subset_completed}

Los resultados CDR cuantifican un sistema base ya completado que contiene el packing distinguido. No deducen la existencia de esa base a partir de solas desigualdades de tamaños de clases. Las familias finales son conjuntos reales de bloques, no copias marcadas. Los cuadrados locales se recuperan de pares fijos fila-columna, por lo que su conteo producto no tiene pérdida de multiplicidad. Los conteos de tipos se refieren a multiconjuntos de colores y coinciden con los de la base; se conservan las configuraciones fijas del residuo, sin afirmar que su número total quede igual.

\subsection{5. Herramientas, imports y reconstrucción}\label{herramientas-imports-y-reconstrucciuxf3n}

El proyecto actual fija Lean 4.35.0, el commit de Mathlib \path{56a005855ea17e0b0958f73ec801ae140b0a1679} y otros ocho repositorios del ecosistema. Declara 310 raíces de módulos locales. La entrada canónica \path{LatinonFormalization} importa el ensamblaje de aplicaciones y \path{CorrectionFive}; su clausura local de imports tiene 299 módulos. El manifiesto de fuentes lista todos los módulos y hashes de archivos seleccionados.

Extraiga \path{preprints/LATINON/05_formalization/Latinon_Lean_v1.0.zip}, que contiene el proyecto \path{lean/} y las licencias de origen. Instale el toolchain y los paquetes externos fijados y ejecute secuencialmente \texttt{lake\ build\ LatinonFormalization} y \texttt{lake\ env\ lean\ ReviewAxioms.lean} desde ese proyecto. El build de la biblioteca abarca 310 objetivos configurados. La reconstrucción original de 309 objetivos comenzó sin objetos del proyecto; el build incremental posterior reutilizó esos objetos y las cachés externas fijadas.

Ambos builds terminaron el 16 de noviembre de 2026. El archivo actual de doce consultas y su salida exitosa se entregan con la evidencia incremental. Todos los conjuntos de axiomas informados están contenidos en \path{propext}, \path{Classical.choice} y \path{Quot.sound}. El archivo y los logs están vinculados mediante SHA-256; la verificación del paquete de solo lectura no ejecuta Lean.

Ejecute \texttt{python\ scripts/review\_axioms.py} desde el repositorio de publicación para verificar las doce consultas registradas con un solo comando. Comprueba los vínculos de logs y fuentes certificados. Para una ejecución nueva de Lean, utilice el comando de consultas anterior.

\subsection{6. Procedencia y límites de la revisión}\label{procedencia-y-luxedmites-de-la-revisiuxf3n}

La instantánea limpia original se identifica mediante \path{7b603942c8d078db6c0a5c8729df7012b7786db8} y se conserva por separado. El freeze actual añade un módulo y modifica la entrada canónica, el archivo de consultas y la configuración de la biblioteca. Su manifiesto de fuentes se regenera para el ZIP entregado. \path{preprints/LATINON/02_validation/v1.0/FREEZE.json} vincula el archivo actual y los hashes de sus entradas; \path{BASELINE_FREEZE.json} vincula la evidencia original. Desde la raíz del repositorio, \texttt{python\ scripts/verify\_repository.py} comprueba ambos sin objetos Git privados.

Nueve módulos de grafones provienen del proyecto de Cameron Freer en el commit \path{9f7be59fa754d260a544b4cfd83d6a5b94f7552e}. El archivo de procedencia conservado enumera los metadatos blueprint retirados y los cambios de compatibilidad. La selección CDR finita proviene de \path{jtraverso/cdr-palomar-kernel} en el commit \path{5d1e6fe301e5fc01c92f6145aa34fddfe9c2bd10}. Las definiciones básicas son idénticas byte a byte; se enumeran los cambios de compatibilidad de Wilson y de representación de matrices latinas. Esta selección no importa la antigua biblioteca de existencia por absorción.

La auditoría interna sigue el protocolo del Paper III de CDR en el commit \path{d12d980f8c86d8795e9faa0c6bb274ee936c7c8e} de la serie. Separa reconstrucción, axiomas, enunciados protegidos, revisión dirigida de pruebas escritas, paridad bilingüe, documentos renderizados e integridad de archivos. La auditoría de la edición es \path{publication/LATINON_RELEASE_AUDIT.json}. Es una comprobación interna de preparación, no revisión independiente por pares ni certificación exhaustiva de prioridad.

\end{document}
